\documentclass[aps,preprint]{revtex4}%
\usepackage{amsfonts}
\usepackage{amsmath}
\usepackage{amssymb}
\usepackage{graphicx}%
\setcounter{MaxMatrixCols}{30}
%TCIDATA{OutputFilter=latex2.dll}
%TCIDATA{Version=4.00.0.2321}
%TCIDATA{CSTFile=revtex4.cst}
%TCIDATA{Created=Thursday, August 24, 2006 16:22:24}
%TCIDATA{LastRevised=Friday, August 25, 2006 12:47:02}
%TCIDATA{<META NAME="GraphicsSave" CONTENT="32">}
%TCIDATA{<META NAME="DocumentShell" CONTENT="Articles\SW\REVTeX 4">}
\newtheorem{theorem}{Theorem}
\newtheorem{acknowledgement}[theorem]{Acknowledgement}
\newtheorem{algorithm}[theorem]{Algorithm}
\newtheorem{axiom}[theorem]{Axiom}
\newtheorem{claim}[theorem]{Claim}
\newtheorem{conclusion}[theorem]{Conclusion}
\newtheorem{condition}[theorem]{Condition}
\newtheorem{conjecture}[theorem]{Conjecture}
\newtheorem{corollary}[theorem]{Corollary}
\newtheorem{criterion}[theorem]{Criterion}
\newtheorem{definition}[theorem]{Definition}
\newtheorem{example}[theorem]{Example}
\newtheorem{exercise}[theorem]{Exercise}
\newtheorem{lemma}[theorem]{Lemma}
\newtheorem{notation}[theorem]{Notation}
\newtheorem{problem}[theorem]{Problem}
\newtheorem{proposition}[theorem]{Proposition}
\newtheorem{remark}[theorem]{Remark}
\newtheorem{solution}[theorem]{Solution}
\newtheorem{summary}[theorem]{Summary}
\newenvironment{proof}[1][Proof]{\noindent\textbf{#1.} }{\ \rule{0.5em}{0.5em}}
\begin{document}
\preprint{HEP/123-qed}
\title[Short title for running header]{Long title that appears on the title page}
\author{First Author}
\affiliation{My Institution}
\author{Second Author}
\affiliation{My Institution}
\author{Third Author}
\affiliation{Other Institution}
\keywords{one two three}
\pacs{PACS number}

\begin{abstract}
Shell document for REV\TeX{} 4.

\end{abstract}
\volumeyear{year}
\volumenumber{number}
\issuenumber{number}
\eid{identifier}
\date[Date text]{date}
\received[Received text]{date}

\revised[Revised text]{date}

\accepted[Accepted text]{date}

\published[Published text]{date}

\startpage{101}
\endpage{102}
\maketitle
\tableofcontents
%

\begin{equation}
D\varphi=\left(  \frac{d}{dt}-\varsigma_{1}I\right)  \left(  \frac{d}%
{dt}-\varsigma_{2}I\right)  \varphi=0;\quad\varsigma_{1},\varsigma_{2}%
\in\mathbb{C}%
\end{equation}
$\Rightarrow$%
\begin{equation}
\varphi\in\ker\left(  \frac{d}{dt}-\varsigma_{1}I\right)  \text{ and/or }%
\ker\left(  \frac{d}{dt}-\varsigma_{2}I\right)
\end{equation}%
\begin{equation}
D=\frac{d^{2}}{dt^{2}}-\left(  \varsigma_{1}+\varsigma_{2}\right)  \frac
{d}{dt}+\varsigma_{1}\varsigma_{2}I
\end{equation}%
\begin{align}
\varphi &  =A\varphi_{1}+B\varphi_{2};\varphi_{1}\in\ker\left(  \frac{d}%
{dt}-\varsigma_{1}I\right)  ,\varphi_{2}\in\ker\left(  \frac{d}{dt}%
-\varsigma_{2}I\right)  \\
&  \Downarrow\\
\varphi_{1}\left(  t\right)   &  =e^{\varsigma_{1}t}+\kappa_{1};\quad
\varphi_{2}\left(  t\right)  =e^{\varsigma_{2}t}+\kappa_{2}%
\end{align}
or one make it an initial value problem
\begin{equation}
\varphi_{1}\left(  t\right)  =\varphi_{1}\left(  0\right)  e^{\varsigma_{1}%
t};\quad\varphi_{2}\left(  t\right)  =\varphi_{2}\left(  0\right)
e^{\varsigma_{2}t}%
\end{equation}%
\begin{equation}
\mathcal{L}\left(  D\right)  =z^{2}-\left(  \varsigma_{1}+\varsigma
_{2}\right)  z+\varsigma_{1}\varsigma_{2}I
\end{equation}%
\begin{equation}
\mathcal{L}\left(  D\right)  \mathcal{L}\left(  \varphi\right)  =0\equiv
\left\{  z^{2}-\left(  \varsigma_{1}+\varsigma_{2}\right)  z+\varsigma
_{1}\varsigma_{2}I\right\}  \varphi\left(  z\right)  =0
\end{equation}%
\begin{align}
\mathcal{L}\left(  \varphi_{1}\right)  \left(  z\right)   &  =\int_{0}%
^{\infty}\left\{  e^{\varsigma_{1}t}+\kappa_{1}\right\}  e^{-zt}dt=\int
_{0}^{\infty}\left\{  e^{\left(  \varsigma_{1}-z\right)  t}+\kappa_{1}%
e^{-zt}\right\}  dt=\left\vert \overset{\infty}{\underset{0}{}}\frac
{1}{\left(  \varsigma_{1}-z\right)  }e^{\left(  \varsigma_{1}-z\right)
t}-\frac{1}{z}\kappa_{1}e^{-zt}\right.  \\
&  =\frac{1}{\left(  z-\varsigma_{1}\right)  }+\frac{1}{z}\kappa_{1}%
\end{align}%
\begin{equation}
\varphi_{1}\left(  t\right)  =\left[  \mathcal{L}^{-1}\left(  \mathcal{L}%
\left(  \varphi_{1}\right)  \right)  \right]  \left(  t\right)  =\int
_{C\left(  z\right)  }\left\{  \frac{1}{\left(  z-\varsigma_{1}\right)
}+\frac{1}{z}\kappa_{1}\right\}  e^{zt}dz=e^{\varsigma_{1}t}+\kappa_{1}%
\end{equation}
For the initial value form
\begin{equation}
\int_{0}^{\infty}\varphi_{1}\left(  0\right)  e^{\varsigma_{1}t}%
e^{-zt}dt=\left\vert \overset{\infty}{\underset{0}{}}\frac{\varphi_{1}\left(
0\right)  }{\left(  \varsigma_{1}-z\right)  }e^{\left(  \varsigma
_{1}-z\right)  t}\right.  =\frac{\varphi_{1}\left(  0\right)  }{\left(
z-\varsigma_{1}\right)  }%
\end{equation}
and
\begin{equation}
\int_{C\left(  z\right)  }\frac{\varphi_{1}\left(  0\right)  }{\left(
z-\varsigma_{1}\right)  }e^{zt}dz=\varphi_{1}\left(  0\right)  e^{\varsigma
_{1}t}%
\end{equation}



\end{document}